\section{The depth--shift continuation problem}\label{sec:path}

\subsection{A sufficient diagonal route}

The end of the program requires a sequence of operator bounds,
not a particular smooth curve in the parameter plane.

\begin{proposition}[Sufficient depth--shift route]\label{prop:diagonal}
Suppose $L_j\to\infty$, $0<\omega_j\downarrow0$, and
$\norm{V_{\omega_j,L_j}}\leq1$ for every $j$.
Then $Q_{0,L}[f]\geq0$ for every $L>0$ and
$f\in C_c^\infty(I_L)$, and hence RH holds.
\end{proposition}
\begin{proof}
Fix $L$ and $f$. For large $j$, zero-extend $f$ to $I_{L_j}$.
Restriction of its output to $I_L$ is $V_{\omega_j,L}f$,
because the causal kernel is the same. Restriction cannot
increase its norm, so
$\norm{V_{\omega_j,L}f}\leq\norm f$.
Divide its nonnegative defect by $2\omega_j$ and take the fixed-test
limit \eqref{eq:central-limit}. Theorem~\ref{thm:weil} completes
the implication.
\end{proof}

No rate coupling $L_j$ to $\omega_j$, uniform positive margin,
or differentiable path is needed for this implication. Every
bound must hold on all inputs at its selected pair.
Finite matrices alone do not supply that assertion.

\subsection{What is meant by a critical path}

Here a critical path means a continuation along which a
construction or rigorous estimate remains effective as depth grows
and shift decreases. It is not an assumed boundary of the true
contraction region. Under RH the positive-shift quotient is inner,
so finite-window contraction holds throughout the parameter range;
the parent Wilson-lines discussion makes this distinction
\cite{ParentWilson}. There is also an unconditional anchor at
$\omega=1/2$ from the known innerness of the quotient
\cite[Section~1.4]{Suzuki2012}. This anchor gives finite-window
contraction, but does not automatically initialize the strict
margins or physical data needed for a chosen continuation.

The earlier adaptive-shift estimate is a different, sufficient
route. Assuming $Q_{0,L}\succeq m_LI>0$, and writing $X_Lf=xf$
on the centered interval, put
\begin{equation}
 \kappa_L=\sup_{f\ne0}
       \left(\frac{Q_{0,L}[X_Lf]}{Q_{0,L}[f]}\right)^{1/2}.
\end{equation}
With a finite proved bound on this form-multiplication norm, the
inherited estimate is
\begin{equation}
 Q_{\omega,L}\succeq
 \cos(2\omega\kappa_L)Q_{0,L},
 \qquad 0\leq\omega\kappa_L\leq\pi/4
 \quad\cite{AdaptivePath}.
\end{equation}
It gives a small-shift schedule after central positivity has
been proved at that depth. It does not establish the central
positivity needed to enlarge the interval. A lower estimate of
$\kappa_L$ from finitely many vectors cannot justify that schedule.

\subsection{Length extension has its own exact constraint}

Translate the interval to $(0,L+h)$ and split it into old and
new pieces. Causality gives
\begin{equation}
 V_{\omega,L+h}=
 \begin{pmatrix}X&0\\Y&Z\end{pmatrix},
 \qquad X=V_{\omega,L},\quad Z=V_{\omega,h}.
\end{equation}
If $X$ and $Z$ are strict contractions in operator norm, set
\[
 E=I-X^*X,\qquad F_{\mathrm{out}}=I-ZZ^*.
\]
The inherited cumulative-storage criterion is
\begin{equation}\label{eq:gluing-main}
 \boxed{\norm{V_{\omega,L+h}}\leq1
 \quad\Longleftrightarrow\quad
 \norm{F_{\mathrm{out}}^{-1/2}YE^{-1/2}}\leq1.}
\end{equation}
It measures cross-interval transfer in the old-input and new-output
defect metrics. A strict inequality gives a strict contraction.
The proof and the simultaneous depth--shift criterion are in
\cite[Section~S9]{Supplement}, following \cite{CumulativePath}.
The diagonal strictness is required for the displayed inverses;
a semidefinite boundary case needs a separate range formulation.

This suggests a precise physical gluing target: derive or bound
the normalized cross map in \eqref{eq:gluing-main}, preserving its
directions and cancellations. Separate worst-case scalar estimates
can lose the information a continuation needs.
If the shift also changes, the old and new defect metrics change
as well; the finite-step Schur test in the companion keeps those
changes. The fixed-$L$ shift ODE alone is not that joint step.

For reflection-compatible coordinates one can use
\[
 P_{\omega,L}=\frac2\omega\Ree\bigl[(I-V_{\omega,L})
                                    (I+V_{\omega,L})^{-1}\bigr].
\]
Here $\Ree A=(A+A^*)/2$ and the inverse exists for the
positive-shift finite-window Volterra operator.
The exact congruence
\begin{equation}
 D_{\omega,L}=\frac\omega2(I+V_{\omega,L}^*)P_{\omega,L}
                                      (I+V_{\omega,L})
\end{equation}
preserves positivity and the normalized spatial coupling.
These are useful coordinates, not an additional positivity result.
At positive shift $P$ is bounded, whereas $Q_{0,L}$ is an unbounded
form; a core expansion does not give a uniform operator-norm
approximation.

\subsection{Separate fixed-window arithmetic progress}

The separate EMA calculation now supplies an internally certified
all-input bound at a designated small window:
\begin{equation}\label{eq:main-ema-anchor}
 Q_{0,1/2}\succeq I/40,\qquad
 D_{10^{-3},1/2}\succeq0.000049998\,I.
\end{equation}
It uses the fixed arithmetic tower, exact endpoint response and
rigorous bounds on omitted inputs, with rational interval checks
at two precisions \cite{EMAAnchor}. Specialist review remains open.
This closes the pilot's input-complement gap, without claiming a
new depth horizon or an all-depth result.

For the adjoining window $h=1/20$ at the same shift, the diagonal
floor transfers to both defect orientations. The scalar sufficient
estimate loses a factor between 55 and 72, and the finite quotient
near $0.802$ is not an all-input upper bound \cite{CumulativeAppend}.
A subsequent energy-weighted operator calculation establishes
internally, at this designated append,
\[
 \norm{F_{\mathrm{out}}^{-1/2}YE^{-1/2}}<0.951<1
 \qquad\cite{CumulativeCoupling}.
\]
Its forward/backward energy argument controls the entire input
complements; specialist review remains outstanding. It closes
this bounded append test, without a new positivity horizon or an
all-depth continuation. Its proof remains in the separate
critical-path investigation. Neither it nor the new physical
hierarchy in Section~\ref{sec:growing-trace} identifies a Wilson
expectation with the arithmetic transfer.

\subsection{The all-depth requirement}

A recursion with $L_{j+1}=L_j+h_j$ must establish
$\sum_jh_j=\infty$ as well as $\omega_j\downarrow0$.
Local extendibility could otherwise accumulate at a finite length.
Arithmetic delays and crossings of weak directions favor a finite
step criterion over an assumed smooth critical-curve ODE.
Decreasing the shift also cannot be justified by reversing forward
dissipation; no bounded inverse of $V_{\omega,L}$ is assumed.

Numerical exploration should retain complete output norms and a
rigorous bound on the unresolved input complement. In particular,
$I-(\Pi V\Pi)^*(\Pi V\Pi)$ can look positive because output was
discarded. The parent Loewner assembly is a reusable diagnostic
tool, but a finite experiment neither proves
\eqref{eq:gluing-main} on the full space nor prevents finite-depth
accumulation.
